%\ifodd\value{page}\hbox{}\newpage\fi
\chapter*{Śrī Lalitā Sahasranāma — Śloka 15}
\addcontentsline{toc}{chapter}{Śloka 15 - Nomi 29,30}

\begin{sloka}
{\sanskritfont
\begin{tabular}{c}
लक्ष्यरोमलताधारतासमुन्नेयमध्यमा ।\\
स्तनभारदलन्मध्यपट्टबन्धवलित्रया ॥ १५॥
\end{tabular}

\vspace{1.5em}

{\middlesize \iastfont
lakṣya-roma-latā-dhāratā-samunneya-madhyamā |\\
stana-bhāra-dalan-madhya-paṭṭa-bandha-valitrayā ||15||}

\vspace{1em}

{\middlesize
\anvaya{%
लक्ष्यरोमलताधारतासमुन्नेयमध्यमा ।
स्तनभारदलन्मध्यपट्टबन्धवलित्रया ।
}
}

\middlesize \iastfont \itmean{%
«Il cui punto mediano è reso percepibile dalla linea dei peli addominali;  
il cui ventre è segnato da tre pieghe, dovute al peso dei seni e alla cintura che lo avvolge.»%
}

}
\end{sloka}

\wordmeaning{%
\wordline{लक्ष्य-रोम-लता-धारता-समुन्नेय-मध्यमा}
{lakṣya-roma-latā-dhāratā-\\samunneya-madhyamā}
{ il cui centro (\textit{madhyamā}) è reso percepibile (\textit{samunneya}) dalla continuità (\textit{dhāratā}) della linea di peli sottili (\textit{roma-latā}), chiaramente distinguibile (\textit{lakṣya}); }%

\wordline{स्तन-भार-दलन्-मध्य-पट्ट-बन्ध-वलि-त्रया}
{stana-bhāra-dalan-madhya-\\paṭṭa-bandha-valitrayā}
{ il cui ventre (\textit{madhya}) presenta tre pieghe (\textit{vali-trayā}), causate dalla pressione (\textit{dalan}) del peso dei seni (\textit{stana-bhāra}) e dalla fascia o cintura che lo avvolge (\textit{paṭṭa-bandha}); }%
}

\textit{Śloka} 15 prosegue la contemplazione del corpo della Dea concentrandosi sulla regione mediana, luogo simbolico di equilibrio e sostegno. La \textit{roma-latā}, la sottile linea di peli che discende dall’ombelico, non è descritta come ornamento sensuale, ma come segno rivelatore: essa rende percepibile il centro, il punto di tenuta del corpo e della presenza.

Il termine \textit{samunneya} indica ciò che non è immediatamente dato, ma che può essere riconosciuto attraverso un segno. Il ventre della Dea non si impone allo sguardo: si lascia comprendere come asse silenzioso, come misura.

Nel secondo \textit{pāda}, le tre pieghe dell’addome non alludono a imperfezione o peso corporeo, ma al naturale adattamento della forma alla sua pienezza. Il peso dei seni e la cintura che avvolge il ventre producono una triplice ondulazione che diviene segno di stabilità, non di tensione.

Questo \textit{śloka} contiene due \textit{nāma}.  
\textbf{\textit{Lakṣya-roma-latā-dhāratā-\\samunneya-madhyamā}} descrive il centro della Dea come riconoscibile attraverso un segno sottile e continuo.  
\textbf{\textit{Stana-bhāra-dalan-madhya-paṭṭa-bandha-valitrayā}} presenta la triplice piega del ventre come espressione naturale dell’armonia tra peso, sostegno e contenimento.
